\chapter{Forecasting and regression}
\label{ch:pronosticos}

\section{A forecast is an operational hypothesis}
Forecasts inform inventory, capacity, and planning. A straight line that fits
the past well does not guarantee future accuracy
\cite{hyndman2021forecasting}. QSB implements linear trend, moving average,
simple exponential smoothing, multiplicative seasonal decomposition, and
least-squares linear regression. The JSON is a request: it contains the model,
method, horizon, and parameters. This preserves the chosen method alongside
the data. Native and validation backends are available; external integration
is a future extension.

\section{Data and errors}
Observations are processed in row order; labels do not reorder dates or fill
missing periods. A monthly series must have consistent monthly spacing. QSB
returns fitted values, residuals \(e_t=y_t-\hat y_t\), MAD, MSE, and MAPE:
\[
 \mathrm{MAD}=\frac1m\sum|e_t|,\quad
 \mathrm{MSE}=\frac1m\sum e_t^2,\quad
 \mathrm{MAPE}=\frac{100}{m_+}\sum_{|y_t|>10^{-12}}|e_t/y_t|.
\]
\(m\) counts the available fitted rows, rather than necessarily the entire
history; \(m_+\) excludes actual values whose magnitude is no greater than
\(10^{-12}\). If all are excluded, MAPE is unavailable. MAD retains the
series' units; MSE has squared units; MAPE is reported as a percentage.
Comparing methods requires considering that moving average and smoothing omit
different initial rows.

\section{Linear trend: a complete example}
The model \(y_t=a+bt+e_t\) estimates \(a,b\) by least squares with indices
\(t=1,\ldots,n\). \path{manual/examples/prob_pronostico.json} contains
\((12,14,16,18,20,22)\), month as its unit, and a horizon of two. From the
root:
\begin{lstlisting}
.build/debug/qsb validate-forecast-json manual/examples/prob_pronostico.json
.build/debug/qsb solve-forecast-json manual/examples/prob_pronostico.json
\end{lstlisting}
QSB obtains \(a=10,b=2\), forecasts of \(24\) and \(26\), and zero for
all three error measures. These artificial data were chosen to check the
formula and indices: zero error is not evidence about real demand. The result
retains the request, the \texttt{linearTrend} algorithm, and \texttt{exact}
exactness, which describes the method's calculation rather than exact knowledge
of the future.

\section{Moving average and smoothing}
For window \(k\), the one-step forecast is
\[
 \hat y_{t+1}=\frac1k\sum_{j=0}^{k-1}y_{t-j}.
\]
QSB extends horizons by using earlier forecasts within the window when actual
values run out. The file \path{prob_movil.json} selects a window of three
for the same series. The results are \(20\) and \(20.666667\): the second
average uses \((20,22,20)\). MAD is \(4\), MSE \(16\), and MAPE
\(20.134680\%\), calculated over the final three rows. This lag is expected
when the series has a trend.

Simple smoothing applies
\[
 F_{t+1}=\alpha y_t+(1-\alpha)F_t,\qquad 0<\alpha\le1.
\]
The initial forecast is set to the first observation. In
\path{prob_suavizado.json}, \(\alpha=0.3\) produces \(18.117660\) in
both future periods. The method has no trend component; it does not
automatically extrapolate growth. Its MAD is \(4.078440\), MSE
\(18.233818\), and MAPE \(22.081810\%\). A large alpha reacts quickly
and may follow more noise; a small one smooths and may lag.
Both examples run with \texttt{solve-forecast-json} and are validated with
\texttt{validate-forecast-json}.

\section{Multiplicative seasonality}
The method combines the estimated trend with seasonal factors:
\(\hat y_t=(a+bt)s_j\). QSB calculates actual/trend ratios, averages them
by seasonal position, and normalizes the factors to have a mean of one.
This is an educational decomposition, rather than Holt--Winters or an ARIMA
model. Seasonal length and observation order must correspond to the phenomenon;
at least two complete cycles and values/trend compatible with multiplicative
division are required.
\path{prob_estacional.json} uses length two for the previous linear series.
Both factors are one; the forecast is again \(24,26\). The case checks that
specifying a season does not invent oscillations when data coincide with the
line. To study actual seasonality, replace the series with several oscillating
cycles and review residuals before interpreting the factors.

\section{Explanatory regression}
The model is \(y_i=\beta_0+\sum_j\beta_jx_{ij}+e_i\), with an intercept.
Predictors must have linearly independent columns; a constant column duplicates
the intercept. The core detects rank deficiency.
\path{prob_regresion.json} contains activity \(1,2,3,4,5\) and consumption
\(5,7,9,11,13\). With \texttt{solve-forecast-json}, QSB returns intercept
\(3\), activity coefficient \(2\), zero SSE, and \(R^2=1\).
The request contract's horizon must also be positive for regression, although
its output fits the supplied observations and does not generate future predictor
rows. For a new activity level, the estimated equation allows a prediction to
be calculated, without QSB reporting an interval for it.

A coefficient is a conditional association, rather than causality. \(R^2\)
measures within-sample fit; it does not identify overfitting or guarantee
extrapolation. This workflow does not report coefficient standard errors,
hypothesis tests, or prediction intervals. Small original datasets allow the
contract to be checked, but do not assess generalization ability.

\section{Working in macOS and evaluation}
In the forecasting editor, select the method, parameters, and horizon; enter
historical rows in temporal order. Values can be pasted into one column, or
labels and tab-separated values into two columns. Paste validation is atomic:
a malformed row does not partially replace the history. Regression uses
variable names and explanatory rows. Validate, solve, and review actual,
fitted, and future curves; these curves and the residual chart share observation
order. Visibility and horizontal-scale controls help with long series. The
summary shows parameters and metrics; the JSON selector preserves the request
and result. Applying JSON updates the draft only when decoding succeeds.

To evaluate an operational forecast, reserve the final periods, fit using the
earlier history, and compare against the held-out data. The errors QSB reports
are the fits from the executed calculation; they do not automatically constitute
an out-of-sample test. Do not select parameters by repeatedly examining the
same held-out set. Document units, window, alpha, season, and horizon when
delivering the analysis.

\section{Exercises}
\begin{enumerate}
\item Compare windows of two and four for the growing series. Explain the lag
and check which rows contribute to MAD.
\item Replace one observation with zero and check the definition of MAPE.
Compare its suitability with MAD.
\item Enter a second predictor equal to twice the first and observe the rank
diagnostic. Explain why more columns do not always add information.
\end{enumerate}
